\message{ !name(main.tex)}\documentclass[a4paper, margin=0.5in]{article}
\usepackage[utf8]{inputenc}
\usepackage[T1]{fontenc}
\usepackage{textcomp}
\usepackage[english]{babel}
\usepackage{amsmath, amssymb, amsthm}
\usepackage{hyperref}
\usepackage{cleveref}
\usepackage{geometry}

% figure support
\usepackage{import}
\usepackage{xifthen}
\pdfminorversion=7
\usepackage{pdfpages}
\usepackage{transparent}
\newcommand{\incfig}[1]{%
  \def\svgwidth{\columnwidth}
  \import{./Figures/}{#1.pdf_tex}
}

\usepackage[most]{tcolorbox}

\pdfsuppresswarningpagegroup=1

\newcommand{\R}{\mathbb{R}}
\newcommand{\Q}{\mathbb{Q}}
\newcommand{\Z}{\mathbb{Z}}
\newcommand{\N}{\mathbb{N}}
\renewcommand{\H}{\mathbb{H}}
\newcommand{\F}{\mathbb{F}}
\newcommand{\W}{\mathbb{W}}

\usepackage{tikz}
\usepackage{xparse}

% Define a new environment for polar grid
\NewDocumentEnvironment{polargrid}{O{4}O{4}O{12}O{4.5}}{%
  % #1 = max radius (default: 4)
  % #2 = number of concentric circles (default: 4)
  % #3 = number of rays (default: 12)
  % #4 = axis length beyond max radius (default: 4.5)
  \begin{tikzpicture}[>=latex]
    % Draw the outer circle
    \draw[lightgray, thin] (0,0) circle (#1);
    % Draw concentric circles
    \foreach \r in {1,...,\numexpr#2-1\relax} {
      \draw[lightgray, thin] (0,0) circle (\r*#1/#2);
    }
    % Draw rays
    \foreach \i in {0,...,\numexpr#3-1\relax} {
      \pgfmathsetmacro{\angle}{(\i*360)/#3}
      \draw[lightgray, thin] (\angle:#1) -- (\angle:0);
    }
    % Draw main axes
    \foreach \angle in {0,90,180,270} {
      \draw[black, thick] (\angle:#1) -- (\angle:0);
    }
    % Angle labels
    \foreach \i/\label in {
      30/$\frac{\pi}{6}$,
      60/$\frac{\pi}{3}$,
      120/$\frac{2\pi}{3}$,
      150/$\frac{5\pi}{6}$,
      180/$\pi$,
      210/$\frac{7\pi}{6}$,
      240/$\frac{4\pi}{3}$,
      270/$\frac{3\pi}{2}$,
      300/$\frac{5\pi}{3}$,
      330/$\frac{11\pi}{6}$
    } {
      \draw (\i:#1 + 0.3) node {\label};
    }
    % Draw axis labels
    \draw[->, thick] (0,0) -- (0,#4) node[above] {$\mathbf{y}$};
    \draw[->, thick] (0,0) -- (#4,0) node[right] {$\mathbf{x}$};
  }{%
  \end{tikzpicture}
}

% Command for plotting points on polar grid
% This version properly handles negative radii by adding 180 degrees to the angle
\NewDocumentCommand{\polarpoint}{mmO{red}O{}}{%
  % #1 = radius
  % #2 = angle in radians (will be converted to degrees)
  % #3 = point color (default: red)
  % #4 = label (optional)

  % Handle negative radius by adding 180° to angle and using absolute radius
  \pgfmathsetmacro{\radius}{abs(#1)}
  \pgfmathsetmacro{\angle}{#2*180/pi}

  % If radius is negative, adjust the angle by adding 180°
  \ifdim #1 pt < 0 pt
    \pgfmathsetmacro{\angle}{\angle + 180}
  \fi

  % Normalize angle to be between 0° and 360°
  \pgfmathsetmacro{\angle}{mod(\angle, 360)}

  \draw[fill=#3] (\angle:\radius) circle (2pt) 
  \ifx\relax#4\relax\else
    node[above right, font=\scriptsize] {#4}
  \fi;
}

% For convenience, a command that formats polar coordinates nicely
\NewDocumentCommand{\polarcoord}{mm}{%
  % #1 = radius
  % #2 = angle (in radians)
  $(#1, \formatangle{#2})$%
}

% Helper command to format angles nicely
\NewDocumentCommand{\formatangle}{m}{%
  \ensuremath{#1}%
}

\usepackage[most]{tcolorbox}
\tcbuselibrary{listings, breakable, skins}

% Problem box
\newtcolorbox[auto counter, number within=section]{problem}[1][]{%
  colback=red!5!white,
  colframe=red!100!black,
  fonttitle=\bfseries,
  title=Problem~\thetcbcounter,
  label=#1,
  enhanced,
  breakable,
  sharp corners,
  boxrule=0.8pt,
  coltitle=black,
  before upper={\parindent15pt}, % optional formatting
}

% Solution box
\newtcolorbox[auto counter, number within=section]{solution}[1][]{%
  colback=green!5!white,
  colframe=green!100!black,
  fonttitle=\bfseries,
  title=Solution~\thetcbcounter,
  label=#1,
  enhanced,
  breakable,
  sharp corners,
  boxrule=0.8pt,
  coltitle=black,
  before upper={\parindent15pt}, % optional formatting
}

\author{Bhashkar Paudyal}
\title{Practice Set 1}

\begin{document}

\message{ !name(main.tex) !offset(-3) }

\maketitle
\newpage

\tableofcontents{}
\newpage

\section{Polar Graphing}
\label{sec:polar-graphing}

\begin{problem}
  Plot the point with the given polar coordinates

  \begin{enumerate}
    \item \((-1,\frac{3\pi}{4})\)
    \item \((-4, -\frac{5\pi}{6})\)
    \item \((-2, -\frac{2\pi}{3})\)
    \item \((-1, -\frac{\pi}{3})\)
  \end{enumerate}

  \begin{solution}
    Plot:\\
    \begin{polargrid}[5][5][12][5.5]
      \polarpoint{-1}{3*pi/4}[red][\((-1,\frac{3\pi}{4})\)]
      \polarpoint{-4}{-5*pi/6}[red][\((-4,-\frac{5\pi}{6})\)]
      \polarpoint{-2}{-2*pi/3}[red][\((-2,-\frac{2\pi}{3})\)]
      \polarpoint{-1}{-pi/3}[red][\((-1,-\frac{\pi}{3})\)]
    \end{polargrid}
  \end{solution}
\end{problem}

\newpage

\begin{problem}
  Find all pairs of polar coordinates that describe the same point as the
  provided polar coordinates.\\
  \begin{polargrid}[4][4][24][4.5]
    \polarpoint{-4}{pi/12}[red][\((-4,\frac{\pi}{12})\)]
    \polarpoint{-4}{-7*pi/6}[red][\((-4,-\frac{7\pi}{6})\)]
    \polarpoint{2}{-3*pi/2}[red][\((2,-\frac{3\pi}{2})\)]
    \polarpoint{2}{23*pi/12}[red][\((2,\frac{23\pi}{12})\)]
  \end{polargrid}

  \begin{solution}
    Since no point overlap, there are no same points.
  \end{solution}

\end{problem}

\newpage

\textbf{Convert each pair of polar coordinates to rectangular coordinates.}\\

\begin{enumerate}
  \begin{problem}
    \item \( \left( 4, -\frac{\pi}{4} \right) \)
      \begin{solution}
        We have,\\
        \(r = 4\) and \(\theta = -\frac{\pi}{4}\) \\
        We know,\\
        \[
          x = r\cos(\theta) = 4 \cos\left(-\frac{\pi}{4}\right) = 4 \cos\left(\frac{\pi}{4}\right) = 4 \cdot \frac{\sqrt{2}}{2} = 2\sqrt{2}
        \]
        \(\therefore\) \(x = 2\sqrt{2}\).

        Similarly,\\
        \[
          y = r\sin(\theta) = 4\sin\left(-\frac{\pi}{4}\right) = -4\sin\left(\frac{\pi}{4}\right) = -4 \cdot \frac{\sqrt{2}}{2} = -2\sqrt{2}
        \]
        \(\therefore\) \(y = -2\sqrt{2}\).

        \textbf{Rectangular coordinates:} \(\left(2\sqrt{2}, -2\sqrt{2}\right)\)
      \end{solution}
    \end{problem}

    \begin{problem}
    \item \( \left( -1, -\frac{7\pi}{6} \right) \)
      \begin{solution}
        We have,\\
        \(r = -1\) and \(\theta = -\frac{7\pi}{6}\) \\
        We know,\\
        \[
          x = r\cos(\theta) = -1 \cos\left(-\frac{7\pi}{6}\right) = -1 \cos\left(\frac{7\pi}{6}\right) = -1 \cdot \left(-\frac{\sqrt{3}}{2}\right) = \frac{\sqrt{3}}{2}
        \]
        \(\therefore\) \(x = \frac{\sqrt{3}}{2}\).

        Similarly,\\
        \[
          y = r\sin(\theta) = -1 \sin\left(-\frac{7\pi}{6}\right) = \sin\left(\frac{7\pi}{6}\right) = -\frac{1}{2}
        \]
        \(\therefore\) \(y = -\frac{1}{2}\).

        \textbf{Rectangular coordinates:} \(\left(\frac{\sqrt{3}}{2}, -\frac{1}{2}\right)\)
      \end{solution}
    \end{problem}

    \begin{problem}
    \item \( \left( -1, -\frac{11\pi}{6} \right) \)
      \begin{solution}
        We have,\\
        \(r = -1\) and \(\theta = -\frac{11\pi}{6}\) \\
        We know,\\
        \[
          x = r\cos(\theta) = -1 \cos\left(-\frac{11\pi}{6}\right) = -1 \cos\left(\frac{11\pi}{6}\right) = -1 \cdot \frac{\sqrt{3}}{2} = -\frac{\sqrt{3}}{2}
        \]
        \(\therefore\) \(x = -\frac{\sqrt{3}}{2}\).

        Similarly,\\
        \[
          y = r\sin(\theta) = -1 \sin\left(-\frac{11\pi}{6}\right) = \sin\left(\frac{11\pi}{6}\right) = -\frac{1}{2}
        \]
        \(\therefore\) \(y = -\frac{1}{2}\).

        \textbf{Rectangular coordinates:} \(\left(-\frac{\sqrt{3}}{2}, -\frac{1}{2}\right)\)
      \end{solution}
    \end{problem}

    \begin{problem}
    \item \( \left( 1, \frac{7\pi}{4} \right) \)
      \begin{solution}
        We have,\\
        \(r = 1\) and \(\theta = \frac{7\pi}{4}\) \\
        We know,\\
        \[
          x = r\cos(\theta) = 1 \cos\left(\frac{7\pi}{4}\right) = \frac{\sqrt{2}}{2}
        \]
        \(\therefore\) \(x = \frac{\sqrt{2}}{2}\).

        Similarly,\\
        \[
          y = r\sin(\theta) = 1 \sin\left(\frac{7\pi}{4}\right) = -\frac{\sqrt{2}}{2}
        \]
        \(\therefore\) \(y = -\frac{\sqrt{2}}{2}\).

        \textbf{Rectangular coordinates:} \(\left(\frac{\sqrt{2}}{2}, -\frac{\sqrt{2}}{2}\right)\)
      \end{solution}
    \end{problem}
\end{enumerate}

\end{document}

\message{ !name(main.tex) !offset(-150) }
